\subsection{Radical and coradical}\label{sga3:XXI:6.3}

Let \(\cR\) be a root datum. Put
\[
 N=\{x\in M\mid(\alpha^\ast,x)=0
       \text{ for every }\alpha^\ast\in R^\ast\},
\qquad
 N^\ast=M^\ast/\bigl(\VR(R^\ast)\cap M^\ast\bigr).
\]

\begin{lemma}\label{sga3:XXI:6.3.1}
Consider the canonical morphisms
\[
 N\longrightarrow M,
 \qquad
 M^\ast\longrightarrow N^\ast.
\]
They are transposes of one another, and \(N^\ast\) is identified with
the dual of \(N\).
\end{lemma}

\begin{proof}
This follows immediately from \ref{sga3:XXI:1.2.5}.
\end{proof}

\begin{definition}\label{sga3:XXI:6.3.2}
The \emph{coradical} of \(\cR\), denoted
\(\operatorname{corad}(\cR)\), is the trivial root datum
\[
 \operatorname{corad}(\cR)=(N,N^\ast,\varnothing,\varnothing).
\]
\end{definition}

If
\[
 \cR^0=(M,M^\ast,\varnothing,\varnothing),
\]
which is a trivial root datum, there is therefore a morphism
\[
 \operatorname{corad}(\cR)\longrightarrow\cR^0.
\]

\begin{definition}\label{sga3:XXI:6.3.3}
The \emph{radical} of \(\cR\), denoted
\(\operatorname{rad}(\cR)\), is the trivial root datum
\[
 \operatorname{rad}(\cR)=\operatorname{corad}(\cR^\ast)^\ast.
\]
\end{definition}

Thus there is a diagram
\begin{center}
\begin{tikzcd}[column sep=huge,row sep=large]
\operatorname{corad}(\cR) \arrow[rr] \arrow[dr]
  && \operatorname{rad}(\cR) \\
& \cR^0 \arrow[ur] &
\end{tikzcd}
\end{center}
whose transpose is the corresponding diagram for \(\cR^\ast\).

\begin{lemma}\label{sga3:XXI:6.3.4}
The canonical morphism
\[
 u:\operatorname{corad}(\cR)\longrightarrow\operatorname{rad}(\cR)
\]
is an isogeny.
\end{lemma}

\begin{definition}\label{sga3:XXI:6.3.5}
Put
\[
 N(\cR)=K(u)
 =
 M\Big/
 \left(
   \bigl(\VR(R)\cap M\bigr)
   +\bigcap_{\alpha\in R}\ker(\alpha^\ast)
 \right).
\]
There is then a canonical pairing
\[
 N(\cR)\times N(\cR^\ast)\longrightarrow\mathbb Q/\mathbb Z.
\]
\end{definition}

\begin{lemma}\label{sga3:XXI:6.3.6}
One has
\[
 \rgred\bigl(\operatorname{rad}(\cR)\bigr)
 =
 \rgred\bigl(\operatorname{corad}(\cR)\bigr)
 =
 \rgred(\cR)-\rgss(\cR),
\]
and the following conditions are equivalent:
\begin{enumerate}
\item[(i)] \(\cR\) is semisimple;
\item[(ii)] \(\operatorname{rad}(\cR)=0\);
\item[(iii)] \(\operatorname{corad}(\cR)=0\).
\end{enumerate}
\end{lemma}
